\section{第9章 多元函数微分法及其应用}
\subsection{多元函数的基本概念}
\begin{exercise}{96}{定义域与边界}
求 $z=\ln(x-y)+\frac{\sqrt y}{\sqrt{4-x^2-y^2}}$ 的定义域，分别列出各项的约束，并求定义域的内部与边界，判断定义域是否为开集。
\end{exercise}
\begin{exercise}{97}{复合根式与对数约束}
求 $z=\arccos(2x)+\frac{\sqrt{y-\sqrt x}}{\ln(1-x^2-y^2)}$ 的定义域。说明对数真数为正与分母非零分别排除了哪些点，并判断定义域是否包含其全部边界点。
\end{exercise}
\begin{exercise}{98}{空间定义域与连通性}
求 $u=\sqrt{4-x^2-y^2-z^2}+\frac1{\sqrt{x^2+y^2+z^2}-1}$ 的定义域及边界，判断定义域的有界性、开闭性与连通性，并说明原点是否属于定义域。
\end{exercise}
\begin{exercise}{99}{函数复合与约分条件}
设 $f(u,v)=\frac{uv}{u^2+v}$。求 $g(x,y)=f(xy,\frac xy)$ 的表达式和自然定义域，并说明约分后的代数式能否直接用于原定义域之外的点。
\end{exercise}
\begin{exercise}{100}{变量代换与退化点}
已知对所有 $x\ne0$、$y\in\mathbb R$ 有 $f(x+y,\frac yx)=x^2-y^2$。求映射 $(x,y)\mapsto(x+y,\frac yx)$ 的像集，并在像集上求 $f(u,v)$；特别讨论 $v=-1$ 时的信息是否充分。
\end{exercise}
\begin{exercise}{101}{对数代换与可去延拓}
已知 $x>0$、$x\ne1$、$y\in\mathbb R$ 时
\[
f(x-y,\ln x)=\left(\frac yx-1\right)\frac{e^x}{e^y\ln(x^x)}.
\]
求 $f(u,v)$ 的表达式、由已知关系确定的定义域及 $f(1,1)$。对固定 $u$，讨论 $v\to0$ 时是否存在有限极限。
\end{exercise}
\begin{exercise}{102}{振荡因子的控制与推广}
用 $\varepsilon$--$\delta$ 定义证明 $\lim_{(x,y)\to(0,0)}(x^2+y^2)\sin\frac1{\sqrt{x^2+y^2}}=0$。若将正弦项换为任意有界函数 $g(x,y)$，原结论是否仍成立？证明。
\end{exercise}
\begin{exercise}{103}{路径判别与趋近尺度}
证明下列极限不存在，并分别比较沿直线和沿抛物线趋近所提供的信息：
\[
\lim_{(x,y)\to(0,0)}\frac{xy}{x^2+y^2},\qquad
\lim_{(x,y)\to(0,0)}\frac{x^2y}{x^4+y^2}.
\]
对第二个函数，另证明沿每条过原点的直线趋近时极限都为零。
\end{exercise}
\begin{exercise}{104}{等价无穷小与定义域延拓}
求下列极限，并说明哪些表达式在原分母为零的非孤立点处可以作连续延拓：
\begin{enumerate}
\item $\lim_{(x,y)\to(0,5)}\frac{\sin(xy)}x$；
\item $\lim_{(x,y)\to(0,0)}\frac{x^2+y^2}{|x|+|y|}$；
\item $\lim_{(x,y)\to(0,0)}\frac{xy}{\sqrt{2-e^{xy}}-1}$。
\end{enumerate}
\end{exercise}
\begin{exercise}{105}{径向极限与混合变量}
求 $\lim_{(x,y)\to(0,0)}\frac{\ln(1+x^2+y^2)}{(1-\cos\sqrt{x^2+y^2})(1+e^{xy})}$，并说明为什么可以分别处理径向因子与含 $xy$ 的因子，不能只检验坐标轴。
\end{exercise}
\begin{exercise}{106}{无穷远处的不同趋近方式}
求 $x,y$ 同时趋于 $+\infty$ 时 $\left(1+\frac1{xy}\right)^{x^2/(x+y)}$ 的极限，并给出控制其对数的估计。再固定 $y=b>0$，只令 $x\to+\infty$ 求极限，说明趋近条件对结果的影响。
\end{exercise}
\begin{exercise}{107}{指数衰减与区域约束}
求 $x,y$ 同时趋于 $+\infty$ 时 $(x^2+y^2)e^{-(x+y)}$ 的极限，并作估计证明。若只知道 $x+y\to+\infty$，分别讨论附加条件 $x,y\ge0$ 与允许 $x,y$ 为任意实数时，能否保证相同结论；不能时构造一条反例路径。
\end{exercise}
\begin{exercise}{108}{幂指型极限与对数估计}
求 $\lim_{(x,y)\to(0,0)}(x^2+y^2)^{x^2y^2}$。用 $x^2y^2\le\frac14(x^2+y^2)^2$ 控制其对数，并说明沿坐标轴趋近与一般路径趋近是否一致。
\end{exercise}
\begin{exercise}{109}{自然定义域上的连续性}
求 $f(x,y)=\frac{\sqrt{4x-y^2}}{\ln(1-x^2-y^2)}$ 的自然定义域，说明它在定义域内的连续性，并求 $\lim_{(x,y)\to(1/2,0)}f(x,y)$。判断根式等号边界上属于定义域的点是否连续。
\end{exercise}
\begin{exercise}{110}{连续延拓与可微性区别}
设 $f(x,y)=\frac{xy}{\sqrt{x^2+y^2}}$（非原点），$f(0,0)=a$。求使 $f$ 在原点连续的 $a$，并用定义判断此时两个偏导数及全微分是否存在。
\end{exercise}
\begin{exercise}{111}{振荡项与连续性}
讨论下列函数在原点的连续性。对连续者给出统一估计，对不连续者给出明确的趋近序列或路径：
\begin{enumerate}
\item $f(0,0)=0$，非原点处 $f(x,y)=(x^2+y^2)\ln(x^2+y^2)$；
\item $g(0,0)=0$，非原点处 $g(x,y)=\frac{x^2\sin\frac1{x^2+y^2}+y^2}{x^2+y^2}$。
\end{enumerate}
\end{exercise}
\begin{exercise}{112}{径向振荡与极限条件}
对实数 $\alpha$，设 $F_\alpha(0,0)=0$，非原点处
\[
F_\alpha(x,y)=(x^2+y^2)^{\alpha/2}\sin\frac1{\sqrt{x^2+y^2}}.
\]
完整讨论原点处的极限、连续性及局部有界性。没有极限时，构造两列趋近原点的点说明原因。
\end{exercise}
\begin{exercise}{113}{分子阶数与二重极限}
设 $p,q>0$，$G_{p,q}(0,0)=0$，非原点处 $G_{p,q}(x,y)=\frac{|x|^p|y|^q}{x^2+y^2}$。求使 $G_{p,q}$ 在原点连续的全部 $(p,q)$，并分别讨论原点邻域内有界的充要条件与二重极限存在的充要条件。
\end{exercise}
\subsection{偏导数}
\begin{exercise}{114}{偏导数定义与局部信息}
设 $f$ 在原点邻域有定义，且 $x\ne0$ 时 $\frac{f(x,0)-f(0,0)}x=2x+3$。用定义求 $f_x(0,0)$，并判断这些条件是否足以确定 $f_y(0,0)$，用两个具体函数说明。
\end{exercise}
\begin{exercise}{115}{根式函数的轴向差商}
设 $f(x,y)=e^{\sqrt{x^2+y^4}}$。用定义讨论原点两个偏导数是否存在，并进一步判断原点的连续性及可微性，给出相应理由。
\end{exercise}
\begin{exercise}{116}{偏导数与变步长差商}
设 $f_x(x_0,y_0)$ 存在。分别求
\[
\lim_{h\to0}\frac{f(x_0,y_0)-f(x_0-h,y_0)}h,\qquad
\lim_{h\to0}\frac{f(x_0-2h,y_0)-f(x_0,y_0)}h.
\]
再证明对任意实数 $a$，$\lim_{h\to0}\frac{f(x_0+ah,y_0)-f(x_0,y_0)}h=a f_x(x_0,y_0)$，并单独说明 $a=0$ 的情形。
\end{exercise}
\begin{exercise}{117}{不对称差商与条件辨析}
设 $f_x(x_0,y_0)$ 存在，求 $\lim_{h\to0}\frac{f(x_0+2h,y_0)-f(x_0-3h,y_0)}h$。反过来，仅由这个极限存在，能否推出 $f_x(x_0,y_0)$ 存在？给出证明或反例。
\end{exercise}
\begin{exercise}{118}{绝对值乘积的充分条件}
设 $f(x,y)=|x-y|g(x,y)$，$g$ 在原点邻域有定义。逐项判断下列条件是否足以保证 $f$ 在原点的两个偏导数存在，证明充分者并为不充分者给出反例：$g(0,0)=0$；$\lim g=0$ 且 $g(0,0)=0$；$\lim g$ 存在；$g$ 在原点连续且取值为零。再设 $\lim g=L$，求偏导数存在的充要条件。
\end{exercise}
\begin{exercise}{119}{轴向奇异与连续性}
设 $f(0,0)=0$，非原点处 $f(x,y)=\frac{\sqrt{|x|}}{x^2+y^2}\sin(x^2+y^2)$。讨论原点处两个偏导数是否存在，并证明函数在原点连续，说明连续性与偏导数存在性在本题中的区别。
\end{exercise}
\begin{exercise}{120}{偏导数与局部变化率}
设 $f(x,y)=x^3+3xy+4y-1$。求 $(1,2)$ 处的两个偏导数，并求使 $f(1+h,2+kh)-f(1,2)=o(h)$ 的常数 $k$。
\end{exercise}
\begin{exercise}{121}{特殊点求导与局部定义}
设 $f(x,y)=x+2y+(y-1)\arcsin\frac{x}{1+xy}$。说明原点坐标平移后的点 $(0,1)$ 有一个邻域包含于自然定义域内部，求 $f_x(0,1)$、$f_y(0,1)$ 及 $f_{xy}(0,1)$。
\end{exercise}
\begin{exercise}{122}{立方根的偏导数定义}
设 $f(x,y)=\sqrt[3]{x^5-y^3}$，立方根取实值。用定义求 $f_x(0,0)$、$f_y(0,0)$，并讨论在曲线 $x^5=y^3$ 上的非原点处，两个偏导数是否存在，说明直接套用链式公式的适用范围。
\end{exercise}
\begin{exercise}{123}{初等偏导数与特殊点}
分别求下列函数的两个偏导数，并完成相应判断：
\begin{enumerate}
\item $z=x+2y-\sqrt{x^2+y^2}$；另用定义讨论原点两个偏导数。
\item $z=e^x\sin(x^2+y)$；判断是否存在两个偏导数同时为零的点。
\end{enumerate}
\end{exercise}
\begin{exercise}{124}{幂函数与括号层次}
完成下列求导，并写明所用的实值定义区域：
\begin{enumerate}
\item 设 $u=(\frac xy)^z$，其中 $x/y>0$，求三个一阶偏导数。
\item 在 $x>y$ 上分别设 $U=\ln(x^2+y^2+z^2)+\arctan((x-y)^z)$、$V=\ln(x^2+y^2+z^2)+(\arctan(x-y))^z$。求 $U,V$ 关于 $x,y,z$ 的偏导数，并比较两者关于 $z$ 的求导过程。
\end{enumerate}
\end{exercise}
\begin{exercise}{125}{变底变指数的求导}
在 $xy^2+y>0$ 的区域上，求 $z=(xy^2+y)^x$ 的两个偏导数，并求 $(1,1)$ 处的全微分。写明对数求导中每个变量的作用。
\end{exercise}
\begin{exercise}{126}{齐次复合函数的求导}
在 $x,y>0$ 时设 $z=(\frac yx)^{x/y}$。求 $z_x(1,2)$、$z_y(1,2)$，并证明整个定义区域上恒有 $xz_x+yz_y=0$，解释该恒等式与同时缩放 $x,y$ 的关系。
\end{exercise}
\begin{exercise}{127}{分段函数与偏导连续性}
设 $f(0,0)=0$，非原点处 $f(x,y)=(x^2+y^2)\ln(x^2+y^2)$。求全平面上的两个偏导数，证明它们在原点连续，并判断 $f_{xx}(0,0)$ 是否存在。
\end{exercise}
\begin{exercise}{128}{由偏导数恢复函数}
设 $x>0$，$f_x(x,y)=x+\frac{y^2}x$，$f(1,y)=\sin y$。求 $f$ 及 $f_y$，并说明若定义域改为 $x\ne0$，仅给上述初值是否能唯一确定负半平面上的函数。
\end{exercise}
\begin{exercise}{129}{截线斜率与切向量}
求曲线 $z=\sqrt{1+x^2+y^2},\ x=1$ 在 $(1,1,\sqrt3)$ 处的切线关于 $y$ 轴的斜率，并给出切线参数方程与一个单位切向量。
\end{exercise}
\begin{exercise}{130}{截线倾角与切线方程}
求曲线 $z=\frac{x^2+y^2}4,\ y=4$ 在 $(2,4,5)$ 处沿 $x$ 增大方向的切线倾角，以及该点的切线方程。说明倾角应在哪个平面内理解。
\end{exercise}
\begin{exercise}{131}{坐标单调性与比较}
设 $f\in C^1(\mathbb R^2)$ 且处处 $f_x>0$、$f_y<0$。在 $x_1<x_2,y_1<y_2$；$x_1<x_2,y_1>y_2$；$x_1>x_2,y_1<y_2$；$x_1>x_2,y_1>y_2$ 四组条件中，哪些必推出 $f(x_1,y_1)<f(x_2,y_2)$？用一元中值定理证明，并说明其余条件为何不能保证该结论。
\end{exercise}
\begin{exercise}{132}{连续性与轴向不可导}
证明 $f(x,y)=\sqrt{x^2+y^2}$ 在原点连续但两个偏导数均不存在。进一步证明 $|f(P)-f(Q)|\le |P-Q|$，说明很强的连续性估计仍不能保证偏导数存在。
\end{exercise}
\begin{exercise}{133}{偏导存在与不连续反例}
设 $f(0,0)=0$，非原点处 $f(x,y)=\frac{2xy}{x^2+y^2}$。求原点两个偏导数并证明函数不连续；再求沿 $y=kx$ 的极限，说明两条坐标轴上的信息为何不足。
\end{exercise}
\subsubsection{高阶偏导数}
\begin{exercise}{134}{二阶偏导与调和性}
设 $z=\arctan\frac{x+y}{x-y}$，$x\ne y$。求 $z_{xx},z_{xy}$，并验证 $z_{xx}+z_{yy}=0$。说明混合偏导数可以交换的局部光滑性依据。
\end{exercise}
\begin{exercise}{135}{高阶求导与积分上限}
完成下列计算，并分别说明交换求导顺序的条件：
\begin{enumerate}
\item $f(x,y)=x\sin\frac yx$，求 $f_{xy}(2,\frac\pi2)$ 和 $f_{yx}(2,\frac\pi2)$。
\item $g(x,y)=\int_0^{xy}\frac{\sin t}{1+t^2}\,\dd t$，求 $g_{xx}(0,2)$，并判断是否可以先把 $x=0$ 代入后再连续求两次 $x$ 偏导。
\end{enumerate}
\end{exercise}
\begin{exercise}{136}{三阶偏导与对称性}
设 $z=\ln(e^x+e^y)$。求 $z_{xxy}$，化简后确定其零点集，并比较交换 $x,y$ 后该三阶偏导数的符号变化。
\end{exercise}
\begin{exercise}{137}{径向函数的微分恒等式}
在 $(x,y,z)\ne(0,0,0)$ 上设 $u=\sqrt{x^2+y^2+z^2}$。证明 $u_x^2+u_y^2+u_z^2=1$ 和 $u_{xx}+u_{yy}+u_{zz}=\frac2u$；再直接求导证明 $u^{-1}$ 的三个纯二阶偏导数之和为零。
\end{exercise}
\begin{exercise}{138}{混合偏导不交换的边界}
设 $f(0,0)=0$，非原点处 $f(x,y)=\frac{x^3y}{x^2+y^2}$。分别从定义求 $f_{xy}(0,0)$、$f_{yx}(0,0)$（约定 $f_{xy}=\partial_y(\partial_x f)$），证明两者不等，并讨论这两个混合偏导数在原点的连续性，指出交换定理缺少的条件。
\end{exercise}
\begin{exercise}{139}{轴向信息与混合偏导}
设 $F(0,0)=0$，非原点处 $F(x,y)=\frac{xy(x^2-y^2)}{x^2+y^2}$。用定义求原点的两个一阶偏导数和两个混合偏导数，讨论混合偏导交换的结论是否成立，并证明 $F$ 在原点连续。
\end{exercise}
\begin{exercise}{140}{参数族的一阶与二阶正则性}
对 $a>0$，设 $H_a(0,0)=0$，非原点处 $H_a(x,y)=(x^2+y^2)^a$。分别求使原点处两个一阶偏导数存在、两个纯二阶偏导数存在的 $a$ 的范围，并讨论两个混合偏导数在原点的值；需要时区分端点情形。
\end{exercise}
\subsection{全微分}
\begin{exercise}{141}{全增量与线性近似误差}
设 $f(x,y)=x^2y+xy^2$。求在 $(2,1)$ 处当 $\Delta x=0.1$、$\Delta y=-0.2$ 时的全增量和全微分，并计算两者之差。再对一般增量 $(h,k)$ 写出精确余项，证明其为 $o(\sqrt{h^2+k^2})$。
\end{exercise}
\begin{exercise}{142}{全微分与复合积分}
求下列函数的全微分，并完成相应说明：
\begin{enumerate}
\item $z=(x^2+y^2)e^{\arctan(y/x)}$，$x\ne0$；求在 $(1,0)$ 处的线性近似。
\item $z=\sin x+\sin(xy)+\int_0^{2x+y}e^{-t^2}\,\dd t$；说明变上限积分贡献的两个系数之间有何比例关系。
\end{enumerate}
\end{exercise}
\begin{exercise}{143}{全微分公式的有效区域}
设 $u=\arcsin\frac z{x+y}$，$x+y>0$。求全微分，并写出公式具有有限系数的开区域；说明 $z=\pm(x+y)$ 上不能直接沿用该公式的原因。
\end{exercise}
\begin{exercise}{144}{隐式微分与非退化条件}
设 $z=z(x,y)$ 由 $z-x-y+xe^z=0$ 确定。求 $\dd z$ 并写明非退化条件，再求经过 $(x,y,z)=(0,0,0)$ 的局部分支在原点的全微分。
\end{exercise}
\begin{exercise}{145}{含模长因子的可微性}
设 $f(x,y)=\frac{2x+3y}{1+xy\sqrt{x^2+y^2}}$。从定义证明它在原点可微并求全微分，给出余项关于 $r=\sqrt{x^2+y^2}$ 的一个阶数估计。
\end{exercise}
\begin{exercise}{146}{二阶渐近与一阶微分}
设 $f$ 在原点连续且 $\lim_{(x,y)\to(0,0)}\frac{f(x,y)-x-2y-1}{x^2+y^2}=1$。求 $f(0,0)$ 和 $\dd f|_{(0,0)}$，并写出含二次主项的渐近展开。若去掉原点连续的条件，原点处可微是否仍必然成立？说明理由。
\end{exercise}
\begin{exercise}{147}{全微分的相容性}
设 $u\in C^2(\mathbb R^2)$，$\dd u=3x^2y\,\dd x+f(x,y)\,\dd y$。求 $f_x$ 并说明能否唯一确定 $f$。若另给 $u(0,y)=y^2$，求 $u$ 与 $f$。
\end{exercise}
\begin{exercise}{148}{由全微分恢复函数}
已知 $\dd f=ye^y\,\dd x+[x(1+y)e^y+2y]\,\dd y$，$f(0,0)=2$。求 $f(x,y)$，并通过重新计算两个偏导数验证全微分的两个系数均得到满足。
\end{exercise}
\begin{exercise}{149}{可微定义中的余项尺度}
设 $\Delta f=f(x_0+h,y_0+k)-f(x_0,y_0)$，$r=\sqrt{h^2+k^2}$。判定下列条件中哪些足以保证原点平移后的点 $(x_0,y_0)$ 处可微；对不充分者给出反例：
\begin{enumerate}
\item $\Delta f\to0$；
\item 两个一阶偏导数在该点存在且连续；
\item 两个偏导数存在，且 $\Delta f-f_xh-f_yk\to0$；
\item 两个偏导数存在，且 $\frac{\Delta f-f_xh-f_yk}r\to0$。
\end{enumerate}
比较第三、四项条件中余项尺度的差别。
\end{exercise}
\begin{exercise}{150}{非线性路径与可微必要条件}
设 $f(0,0)=0$，非原点处 $f(x,y)=\frac{3x^2y}{x^4+y^2}$。求原点两个偏导数，并沿适当抛物线判断连续性，从而判定原点可微性。
\end{exercise}
\begin{exercise}{151}{连续但不可微的根式函数}
设 $f(x,y)=\sqrt{x^2+y^4}$。讨论原点处的连续性、两个偏导数和可微性；再说明若仅沿 $y$ 轴观察，为何会遗漏决定性信息。
\end{exercise}
\begin{exercise}{152}{振荡乘积的可微性}
设 $f(0,0)=0$，非原点处 $f(x,y)=xy\sin\frac1{\sqrt{x^2+y^2}}$。用可微定义证明它在原点可微并求全微分；再讨论两个偏导数在原点是否连续，给出明确路径或估计。
\end{exercise}
\begin{exercise}{153}{四类正则性质的逻辑关系}
对某点处的四种性质“连续”“两个偏导数连续”“可微”“两个偏导数存在”，说明以下四条推理链是否正确，并为每条错误链指出失效的箭头，给出相应反例：偏导数连续 $\Rightarrow$ 可微 $\Rightarrow$ 连续；可微 $\Rightarrow$ 偏导数连续 $\Rightarrow$ 连续；可微 $\Rightarrow$ 偏导数存在 $\Rightarrow$ 连续；可微 $\Rightarrow$ 连续 $\Rightarrow$ 偏导数存在。
\end{exercise}
\begin{exercise}{154}{连续与一阶余项检验}
设 $f(0,0)=0$，非原点处 $f(x,y)=\frac{x^2y}{x^2+y^2}$。证明原点连续且两个偏导数存在，再用可微定义检验是否存在全微分；写出使归一化余项不趋于零的具体路径。
\end{exercise}
\begin{exercise}{155}{四种性质的完整辨析}
设 $f(0,0)=0$，非原点处 $f(x,y)=\frac{xy}{\sqrt{x^2+y^2}}$。分别讨论原点的连续性、两个偏导数存在性、偏导数连续性及可微性。计算非原点偏导数，并用它们在坐标轴上的取值核验连续性判断。
\end{exercise}
\begin{exercise}{156}{可微但偏导数不连续}
设 $f(0,0)=0$，非原点处 $f(x,y)=(x^2+y^2)\sin\frac1{x^2+y^2}$。证明它在原点可微并求全微分，证明两个偏导数在原点均不连续；进一步判断它们在原点附近是否有界。
\end{exercise}
\begin{exercise}{157}{余项估计与充分条件的临界阶}
设 $f(0,0)=a$，且原点附近满足 $|f(x,y)-a-bx-cy|\le C(x^2+y^2)^{(1+\alpha)/2}$，其中 $C>0$、$\alpha>0$。证明 $f$ 在原点可微并求其全微分。若将 $\alpha>0$ 放宽为 $\alpha=0$，构造满足估计但不可微的函数。
\end{exercise}
\begin{exercise}{158}{乘积幂函数的可微阈值}
设 $a>0$，$F_a(x,y)=|xy|^a$。讨论原点处的连续性、偏导数及可微性，并求全微分存在的全部参数。再求使 $F_a$ 在原点的某个开邻域内具有连续一阶偏导数的参数范围，说明这与仅在原点可微的条件有何区别。
\end{exercise}
\subsection{多元复合函数的求导法则}
\begin{exercise}{159}{一元复合函数的依赖关系}
设 $z=u^2+v^2+\sin(uv)$，$u=e^x$、$v=x^2$。用链式法则求 $\frac{\dd z}{\dd x}$ 及其在 $x=0$ 处的值，再将 $u,v$ 直接代入求导，核验两种方法一致。
\end{exercise}
\begin{exercise}{160}{全导数与偏导数的区别}
设 $f(x,y)=\arcsin(xy)$，$z(x)=f(x,\ln x)$。在实值可导的开区间内求 $z'(x)$，并分别计算 $z'(1)$ 和 $f_x(1,0)$，说明中间变量的变化怎样进入全导数。
\end{exercise}
\begin{exercise}{161}{退化内映射与一阶导数}
设 $z=f(x\varphi(x),x^2+\varphi(x)^2)$，其中 $f$ 可微，$\varphi$ 可导。求 $z'(x)$，明确 $f$ 的偏导数取值点。若 $\varphi(0)=0$，证明 $z'(0)=0$，无需给出 $f$ 的具体表达式。
\end{exercise}
\begin{exercise}{162}{二阶链式法则}
设 $f\in C^2(\mathbb R^2)$，$z=f(x^2,\sin2x)$。求 $z''(x)$。若 $f_1(0,0)=a$、$f_2(0,0)=b$、$f_{22}(0,0)=c$，求 $z''(0)$，说明哪些数据实际参与结果。
\end{exercise}
\begin{exercise}{163}{内函数驻点与二阶信息}
设 $f\in C^2$ 且 $\dd f|_{(1,1)}=3\,\dd u+4\,\dd v$，$z=f(\cos x,1+x^2)$。求 $z'(0)$、$z''(0)$，并解释为什么不必知道 $f$ 的二阶偏导数的具体值。
\end{exercise}
\begin{exercise}{164}{角度比值与缩放关系}
设 $z=e^u\sin v$，$u=x-y$、$v=\frac yx$，$x\ne0$。求两个一阶偏导数，并证明 $xz_x+yz_y=(x-y)z$。
\end{exercise}
\begin{exercise}{165}{径向复合与原点延拓}
设 $u=\sqrt{x^2+y^2}$，$z=\ln(\arctan u)$。求原点外的两个偏导数及 $xz_x+yz_y$，并判断 $z$ 能否在原点作有限的连续延拓。
\end{exercise}
\begin{exercise}{166}{显含变量与中间变量}
设 $z=e^{x^2+y^2+u^2}$，$u=x^2\sin y$。求 $z_x,z_y$，写出显含 $x,y$ 与通过 $u$ 间接依赖的全部项。再求限制在 $y=0$ 上的一元函数的导数，核验与 $z_x(x,0)$ 一致。
\end{exercise}
\begin{exercise}{167}{链式法则与雅可比退化}
设 $z=f(x\arcsin y,\frac yx)$，$f\in C^1$，$x\ne0$、$|y|<1$。求 $z_x,z_y$，并求 $(1,0)$ 处两个偏导数；说明此点的 $z_x$ 是否与 $f$ 的选择有关。
\end{exercise}
\begin{exercise}{168}{二阶偏导与变量重新表示}
设 $u=x+y$、$v=x^2-y^2$，$z=f(u,v)$，其中 $f\in C^2$。求 $z_{xx},z_{yy}$，并将 $z_{xx}-z_{yy}$ 的结果仅用 $u,v$ 及 $f$ 的偏导数表示。
\end{exercise}
\begin{exercise}{169}{混合偏导的交叉项}
设 $z=f(e^{xy},x^2+y^2)$，$f\in C^2$。求 $z_{xy}$，并化简其在原点的值，说明为什么该点的结果只涉及 $f$ 的一个一阶偏导数。
\end{exercise}
\begin{exercise}{170}{由截面条件恢复外函数}
设 $z=e^{x-y}+f(x-2y)$，且对所有实数 $x$ 有 $z(x,0)=x^3$。先求 $f$，再求 $z_x,z_y,z_{xy}$，并验证恢复出的函数满足给定截面条件。
\end{exercise}
\begin{exercise}{171}{特殊点全微分与零变化方向}
设 $z=yf(x^2-y^2)$，$f$ 可导且 $f(0)=1$、$f'(0)=\frac12$。求 $\dd z|_{(1,1)}$，并求在该点使一阶变化为零的全部非零位移方向。
\end{exercise}
\begin{exercise}{172}{变上下限积分的混合求导}
设 $z=f(2x+y)+f(x-2y)+\int_{x-y}^{x^2+y^2}g(t)\,\dd t$，其中 $f$ 二阶可导，$g$ 一阶可导。求 $z_{xy}$ 及 $z_{xy}(0,0)$，说明上下限两项的符号如何确定。
\end{exercise}
\begin{exercise}{173}{齐次分量与微分恒等式}
设 $z=xy+xf(\frac yx)$，$x\ne0$，$f$ 二阶可导。证明 $xz_x+yz_y=xy+z$ 及 $x^2z_{xx}=y^2z_{yy}$。将 $z$ 分为一次和二次齐次分量，说明第一个恒等式为何不等于通常的一次齐次公式。
\end{exercise}
\begin{exercise}{174}{偏微分约束化为常微分方程}
在 $x,y>0$ 上设 $z=xf(u)$、$u=\frac xy$，$f\in C^1(0,+\infty)$，满足 $xz_x-yz_y=0$ 及 $f(1)=1$。求 $f(u)$，核验所得函数满足约束，并讨论它在 $u=0$ 能否作有限连续延拓。
\end{exercise}
\begin{exercise}{175}{多层依赖的二阶求导}
设 $z=f(2x-y)+g(x,xy)$，$f$ 二阶可导，$g\in C^2$。求 $z_{xx},z_{xy}$；再令 $g(u,v)=uv$，用直接展开的方法核验与 $g$ 有关的各项。
\end{exercise}
\begin{exercise}{176}{三元外函数与二阶链式法则}
设 $u=f(x,y,xe^y)$，$f\in C^2(\mathbb R^3)$。求 $u_{xx},u_{yy}$，并化简 $u_{yy}(0,y)$；说明外函数关于第三变量的一阶偏导项在何种二次求导中出现。
\end{exercise}
\begin{exercise}{177}{嵌套中间变量}
设 $u=f(x,y,z)$，$y=\varphi(x,t)$、$t=\psi(x,z)$，各函数可微，以 $x,z$ 为独立变量。求 $u_x,u_z$。再令 $\varphi(x,t)=x+t$、$\psi(x,z)=xz$，直接消去中间变量检验所得公式。
\end{exercise}
\begin{exercise}{178}{线性变换与二阶算子}
设 $f(x,y)=g(2x+y,3x-y)$，$g\in C^2(\mathbb R^2)$，且 $f_{xx}+f_{xy}-6f_{yy}=1$。求 $g_{uv}$。若另有 $g(u,0)=g(0,v)=0$，求 $g$ 并验证。
\end{exercise}
\begin{exercise}{179}{欧拉公式及二阶推广}
设 $f\in C^2(\mathbb R^2)$，对任意实数 $t$ 有 $f(tx,ty)=t^2f(x,y)$。证明 $xf_x+yf_y=2f$ 及 $x^2f_{xx}+2xyf_{xy}+y^2f_{yy}=2f$；写明两次对参数 $t$ 求导后在哪一点取值，并单独核验原点。
\end{exercise}
\begin{exercise}{180}{全微分形式不变性}
设 $z=\arctan(u+v)$，$u=2x-y^2$、$v=x^2y$。用全微分形式不变性求 $\dd z$，再求 $(1,0)$ 处的全微分，并通过直接代入后的偏导数核验。
\end{exercise}
\begin{exercise}{181}{复合全微分与切向位移}
设 $z=f(2x+3y,e^{xy})$，$f$ 可微。求全微分，并求在 $(0,0)$ 处使 $\dd z=0$ 对每个可微外函数 $f$ 都成立的全部非零位移方向。
\end{exercise}
\begin{exercise}{182}{链式法则所需的外函数条件}
设 $F(0,0)=0$，非原点处 $F(u,v)=\frac{uv}{\sqrt{u^2+v^2}}$，并令 $u(t)=v(t)=t$。证明 $F$ 在原点连续且两个偏导数存在，但 $F(u(t),v(t))$ 在 $t=0$ 不可导。说明只知道外函数偏导数存在时，链式法则的哪项充分条件没有满足。
\end{exercise}
\begin{exercise}{183}{极坐标复合与梯度模}
在 $x>0$ 上设 $r=\sqrt{x^2+y^2}$、$\theta=\arctan\frac yx$，$z=F(r,\theta)$，$F\in C^1$。证明 $z_x^2+z_y^2=F_r^2+\frac1{r^2}F_\theta^2$。若 $F=r^a\cos(b\theta)$、$a,b>0$，求使 $z_x^2+z_y^2$ 在整个右半平面上为常数的全部参数，并求该常数。
\end{exercise}
\subsection{隐函数的求导法则}
\begin{exercise}{184}{隐函数导数与退化点}
由 $e^x+\cos y=xy^2$ 确定的可微分支 $y=y(x)$，求 $y'$ 并写明分母非零条件。核验 $(0,\pi)$ 满足方程，并说明在该点能否直接使用这一求导公式及隐函数存在定理。
\end{exercise}
\begin{exercise}{185}{隐函数初值与切线}
设 $2^{xy}=x+y$ 在 $x=0$ 附近确定 $y=y(x)$。先求 $y(0)$ 并验证隐函数定理的条件，再求 $y'(0)$ 及曲线在该点的切线。
\end{exercise}
\begin{exercise}{186}{二阶隐式求导与常值分支}
设 $y-xe^y+x=0$。求经过 $x=0$ 的局部分支的 $y(0),y'(0),y''(0)$，并用隐函数定理与直接代入说明该局部分支是否恒为常数。
\end{exercise}
\begin{exercise}{187}{含未知函数的隐式求导}
设 $f\in C^1$，$y>0$，$xf(x+y)=e^x\ln y$ 在非退化点确定 $y=y(x)$。求 $y'$ 并写明非退化条件；再求经过 $x=0$ 的分支的初值和 $y'(0)$，结果可含 $f(1)$。
\end{exercise}
\begin{exercise}{188}{隐函数与外函数的二阶复合}
设 $z=f(\ln y-\sin x)$，$f$ 二阶可导且 $f'(0)=1$，$y=y(x)$ 由 $y-xe^{y-1}=1$ 在 $x=0$ 附近确定。求 $z'(0),z''(0)$，并解释为什么无需知道 $f''(0)$ 的值。
\end{exercise}
\begin{exercise}{189}{不同因变量的隐函数条件}
对方程 $xy-z\ln y+e^{xz}=1$，在 $(0,1,1)$ 处计算三个一阶偏导数，判定隐函数定理分别能否保证 $x=x(y,z)$、$y=y(x,z)$、$z=z(x,y)$。再令 $x=0,y=1$，说明本题为何不能唯一确定最后一种局部分支。
\end{exercise}
\begin{exercise}{190}{隐式复合中的结构性抵消}
设 $z+xe^{z-x-y}=x+y$ 在 $1+xe^{z-x-y}\ne0$ 的点附近确定 $z=z(x,y)$。求 $z_x,z_y$，并通过代换 $w=z-x-y$ 解释其中一个偏导数恒为常数的原因。
\end{exercise}
\begin{exercise}{191}{隐函数偏导及恒等式}
设 $xy=xf(z)+y\varphi(z)$，其中 $f,\varphi\in C^1$，$D=xf'(z)+y\varphi'(z)\ne0$。求 $z_x,z_y$，并证明 $xz_x+yz_y=\frac{xy}D$。
\end{exercise}
\begin{exercise}{192}{共同平移下的不变性}
在 $\tan z$ 有定义的局部分支上设 $(x-y)^3+z+\tan z=0$。求 $z_x+z_y$，说明隐函数定理的非退化条件自动成立，并证明在同一局部分支内同时平移 $x,y$ 相同的微小量时，$z$ 保持不变。
\end{exercise}
\begin{exercise}{193}{隐式全微分的非退化条件}
设 $e^z-xy^2+\sin(y+z)=0$ 在某点确定 $z=z(x,y)$。求 $\dd z$，明确给出保证该公式成立的偏导数条件，并写出 $z_x,z_y$；说明为何必须先检查该点满足原方程。
\end{exercise}
\begin{exercise}{194}{隐式微分与图形切平面}
设 $xyz+\sqrt{x^2+y^2+z^2}=\sqrt2$。证明在 $(1,0,1)$ 附近可确定 $z=z(x,y)$，求 $\dd z|_{(1,0)}$，并据此写出该分支图形在此处的切平面。
\end{exercise}
\begin{exercise}{195}{隐函数的二阶偏导}
设 $z^3-3xyz=1$。在 $z^2-xy\ne0$ 的局部分支上求 $z_{xx}$，可保留 $z,z_x$ 的中间表达式，但最终须消去 $z_x$。同时说明方程上的点是否自动满足这一非退化条件。
\end{exercise}
\begin{exercise}{196}{指定分支的二阶偏导和}
设 $z+e^x-y\ln(1+z^2)=0$。求 $(x,y)=(0,0)$ 对应的 $z$ 值，验证局部隐函数条件，并计算 $z_{xx}(0,0)+z_{yy}(0,0)$。
\end{exercise}
\begin{exercise}{197}{直接偏导与复合偏导}
设 $f(x,y,z)=yz^2e^x$，$z=z(x,y)$ 由 $x+y+z+xyz=0$ 确定，并令 $F(x,y)=f(x,y,z(x,y))$。求 $F_x(0,1)$，再分别求 $F_y(0,1)$ 与 $f_y(0,1,-1)$，说明两个记号的含义。
\end{exercise}
\begin{exercise}{198}{隐式中间变量的全微分}
设 $u=f(x,y,z)$，$f\in C^1$，其中 $z=z(x,y)$ 由 $xe^x-ye^y=ze^z$ 确定。求 $\dd u$ 并写明适用的非退化条件，明确区分外函数偏导数与复合函数的全微分系数。
\end{exercise}
\begin{exercise}{199}{低阶数据确定特殊二阶导数}
设 $f\in C^1$、$f(1,1)=0$，$z=z(x,y)$ 由 $(x+1)z-y^2=xf(x+z,x+y)$ 在 $(0,1,1)$ 附近确定。求 $\dd z|_{(0,1)}$、$z_{yy}(0,1)$，并说明求第二个量时为何不必假设 $f$ 有二阶偏导数。
\end{exercise}
\subsubsection{方程组确定的隐函数}
\begin{exercise}{200}{方程组的导数与投影退化}
设 $x+2y+3z=0$、$x^2+y^2+z^2=2$。在能以 $x$ 为参数的局部分支上求 $y',z'$；求曲线上使关于 $(y,z)$ 的雅可比行列式为零的全部点，并说明导数公式的适用范围。
\end{exercise}
\begin{exercise}{201}{隐式方程组与四个偏导}
设 $u^2+xv-y=0$、$yu+v^2-x=0$ 在某点确定 $u=u(x,y),v=v(x,y)$。写出雅可比非退化条件，求 $u_x,v_x,u_y,v_y$，并说明为什么不能不加检查就在 $(x,y,u,v)=(0,0,0,0)$ 使用所求分式。
\end{exercise}
\begin{exercise}{202}{方程组在特殊初值的简化}
设 $f\in C^1$、$f(-1)=1$，$y=y(x),z=z(x)$ 由 $z=xf(x-y)$、$x+y+ze^z=1$ 在 $x=0$ 附近确定。先求初值并核验雅可比条件，再求 $y'(0),z'(0)$，说明为何不需要 $f'(-1)$ 的具体值。
\end{exercise}
\begin{exercise}{203}{两个隐式中间函数的全导数}
设 $u=f(x,y,z)$，$f\in C^1$，$y=y(x)$、$z=z(x)$ 分别由 $e^{xy}+2y=0$ 与 $e^z-xz=2$ 确定。求各非退化分支上的 $\frac{\dd u}{\dd x}$；再求 $x=0$ 处的值，结果以 $f$ 在对应三元点的偏导数表示。
\end{exercise}
\begin{exercise}{204}{非退化条件的充分性}
分别考察方程 $(z-x-y)^3=0$ 与 $z^2-x^2=0$ 在原点附近能否唯一确定连续可微分支 $z=z(x,y)$。计算各方程对 $z$ 的偏导数，并用这两个例子说明隐函数定理中的偏导非零条件与分支存在、唯一之间的逻辑关系。
\end{exercise}
\begin{exercise}{205}{隐函数分支的临界参数}
设 $z+\lambda\sin z=x+y$，考虑经过 $(0,0,0)$ 的分支。求使隐函数定理适用的参数，并在这些参数下求原点的一阶、二阶偏导数。对 $\lambda=-1$，证明仍有唯一连续局部分支，再判断其原点可微性。
\end{exercise}
\subsection{多元函数微分法在几何上的应用}
\begin{exercise}{206}{参数曲线的切线与法平面}
设空间曲线为 $x=t+1$、$y=t^2-1$、$z=t^3$。求 $(2,0,1)$ 处的切线及法平面，写出对应参数与沿参数增大方向的单位切向量，并验证切线方向垂直于法平面。
\end{exercise}
\begin{exercise}{207}{参数法与交面法求切线}
求曲线 $y=x,\ z=x^2-2y$ 在原点的切线与法平面。分别以 $x$ 为参数和用两个曲面法向量的向量积求切向量，说明两种结果一致。
\end{exercise}
\begin{exercise}{208}{平行切线的数量与定位}
对曲线 $y=-x^2,\ z=x^3$，求所有与平面 $x+2y+z=4$ 平行的切线，给出切点及切线方程，并判断这些切线是否可能位于该平面内。
\end{exercise}
\begin{exercise}{209}{交曲线的正则性与法平面}
求曲线 $x^2+y^2+z^2=6,\ x+y+z=0$ 在 $(1,-2,1)$ 处的切线和法平面。先核验该点在曲线上，再证明两个曲面的法向量在此处不平行。
\end{exercise}
\begin{exercise}{210}{隐式曲面的切平面与法线}
求曲面 $x^2+\cos(xy)+yz+x=0$ 在 $(0,1,-1)$ 处的切平面，并求法线参数方程。指出该点附近可由隐函数定理选择哪些坐标作为因变量。
\end{exercise}
\begin{exercise}{211}{椭球切平面与轴截距}
求椭球面 $x^2+2y^2+3z^2=9$ 在 $(2,1,1)$ 处的切平面和法线，进一步求切平面的三个有向轴截距及其到原点的距离。
\end{exercise}
\begin{exercise}{212}{切平面存在与图形表示}
求曲面 $e^z-z+xy=3$ 在 $(2,1,0)$ 处的切平面和法线。说明虽然该点处 $F_z=0$，曲面仍有切平面的原因；再判断隐函数定理能否在此直接保证 $z=z(x,y)$。
\end{exercise}
\begin{exercise}{213}{非标准显式曲面的法向量}
求曲面 $y=x^2+3z^2-15$ 在 $(1,-2,2)$ 处的切平面和法线，并求法线与 $xOz$ 平面的夹角。说明法向量应由哪个隐式函数的梯度得到。
\end{exercise}
\begin{exercise}{214}{平行切平面的反求问题}
求抛物面 $z=4-x^2-y^2$ 上使切平面平行于 $2x+2y+z-1=0$ 的全部点、对应切平面及两个平面的距离，并说明切点是否唯一。
\end{exercise}
\begin{exercise}{215}{切平面包含整条直线}
设 $L:x+y+b=0,\ x+ay-z-3=0$ 位于曲面 $z=x^2+y^2$ 在 $(1,-2,5)$ 处的切平面内。求 $a,b$，并将直线参数式代入切平面，核验包含关系对全部参数成立。
\end{exercise}
\begin{exercise}{216}{零梯度与切平面判定}
平面 $z=0$ 也可写为 $z^3=0$。比较这两种定义函数在原点的梯度，并说明为什么梯度为零不能直接推出曲面没有切平面。再考察圆锥面 $z^2=x^2+y^2$ 在原点的各条母线方向，说明为何不存在能同时承载其一阶方向的单一平面。
\end{exercise}
\begin{exercise}{217}{椭球的指定法向切平面}
设 $\lambda>0$。在椭球面 $x^2+y^2+\lambda z^2=1$ 上求全部点，使其切平面平行于 $x+2y+z=0$；求对应切平面及它们之间的距离，并求切平面到原点的距离等于 $1$ 时的参数 $\lambda$。
\end{exercise}
\subsection{方向导数和梯度}
\begin{exercise}{218}{方向单位化与最大变化率}
求 $f(x,y)=3x^2y-y^2$ 在 $A(1,1)$ 处沿 $A$ 指向 $B(2,1+\sqrt3)$ 的方向导数，写出单位方向向量；再求该点方向导数的最大值并判断给定方向是否达到它。
\end{exercise}
\begin{exercise}{219}{沿曲线正向的方向导数}
设曲线 $y=x^2-5$ 的正向为 $x$ 增大的方向。求 $f(x,y)=x^2-xy+y^2$ 在 $(2,-1)$ 沿曲线切线正向的方向导数，并求反向的结果，说明归一化因子的来源。
\end{exercise}
\begin{exercise}{220}{方向角的相容性}
设 $f(x,y,z)=\ln(x+\sqrt{y^2+z^2})$。先验证 $\alpha=60^\circ,\beta=45^\circ,\gamma=60^\circ$ 可以作为一个空间方向的方向角，再求 $(1,0,1)$ 沿该方向的方向导数及该点方向导数的最大值。
\end{exercise}
\begin{exercise}{221}{单侧方向导数与可微性}
设 $f(x,y)=\sqrt{x^2+y^2}$。采用 $D_{\vec e}f(0)=\lim_{t\to0^+}\frac{f(t\vec e)-f(0)}t$ 的定义，对每个单位方向求方向导数，并讨论连续性、偏导数和可微性。若改用 $t\to0$ 的双侧极限，结论如何？
\end{exercise}
\begin{exercise}{222}{曲面外法向与方向导数}
设 $\vec n$ 为曲面 $2x^2+3y^2+z^2=6$ 在 $P(1,1,1)$ 处的外法向量，$u=\frac{\sqrt{6x^2+8y^2}}z$。求 $P$ 处的外单位法向量及沿该方向的方向导数，再给出沿内法向的结果。
\end{exercise}
\begin{exercise}{223}{梯度与等值线切向}
设 $f(x,y)=\arctan\frac xy$。求 $\nabla f(1,1)$ 及该点使方向导数为零的全部单位方向，说明它们与经过该点的等值线有什么关系。
\end{exercise}
\begin{exercise}{224}{最快增减与零变化方向}
设 $f(x,y,z)=x^3-xy^2-z$，$P(1,1,0)$。求该点增长最快的单位方向及变化率，再求下降最快的单位方向和全部零变化率单位方向，并说明后者形成的几何集合。
\end{exercise}
\begin{exercise}{225}{给定方向与梯度方向比较}
求 $f(x,y,z)=x^2+y^2+xz$ 在 $P(1,0,1)$ 处的梯度，以及沿 $\vec l=(3,-1,6)$ 的方向导数。判断该方向是否为最快增长方向，并给出判断依据。
\end{exercise}
\begin{exercise}{226}{沿另一函数梯度方向求导}
设 $u=\sqrt{x^2+y^2+z^2}$，$g=\frac12x^2+yz$。求在 $(1,1,-1)$ 处沿 $g$ 的梯度方向的 $u$ 的方向导数，判断该方向上 $u$ 是增加还是减少，并与 $u$ 自身的最快增加方向比较。
\end{exercise}
\begin{exercise}{227}{隐函数梯度与方向极值}
设 $x+y^2+z^3=3$ 在 $(x,y)=(1,1)$ 附近确定实值分支 $z=z(x,y)$。验证隐函数条件，求该点方向导数的最大值、最小值及分别取得它们的单位方向。
\end{exercise}
\begin{exercise}{228}{全方向导数存在仍需统一余项}
设 $r=\sqrt{x^2+y^2}$，$F(0,0)=0$，非原点处 $F(x,y)=r\frac{x^3y}{x^6+y^2}$。证明 $F$ 在原点连续且沿任意单位方向的方向导数均为零，再沿 $y=x^3$ 检验可微性。说明所有方向导数均由同一个线性式给出，为何仍不足以保证可微。
\end{exercise}
\begin{exercise}{229}{受切向约束的最速变化}
设温度场 $T(x,y,z)=x^2+2y^2+3z^2$，只能从 $P(1,1,1)$ 沿球面 $x^2+y^2+z^2=3$ 的切向单位方向瞬时移动。求允许方向中温度方向导数的最大值、最小值及取得它们的方向，并求全部零变化率允许方向。用梯度在切平面上的正交投影解释结果。
\end{exercise}
\subsection{多元函数的极值及其应用}
\begin{exercise}{230}{局部极小的必要条件}
设 $f$ 在开区域 $D$ 内可微，内点 $P$ 为局部极小点。证明 $f_x(P)=f_y(P)=0$，并判断能否推出 $f(Q)>f(P)$ 对所有 $Q\in D\setminus\{P\}$ 成立；不能时给出具体反例。
\end{exercise}
\begin{exercise}{231}{极值与坐标截线的导数}
设 $f$ 在内点 $(x_0,y_0)$ 的两个偏导数存在，且在该点取得局部极大值。仅利用两条坐标截线的一元极值必要条件证明两个偏导数为零，并举例说明逆命题不成立。
\end{exercise}
\begin{exercise}{232}{二重渐近确定极值}
设 $f$ 在原点连续，且 $\lim_{(x,y)\to(0,0)}\frac{f(x,y)}{\sin(x^2+y^2)}=1$。求 $f(0,0)$ 和两个偏导数，判断原点是否为严格极值点，并用统一估计证明其类型。
\end{exercise}
\begin{exercise}{233}{四阶主项的符号控制}
设 $f$ 在原点邻域连续，且 $\lim_{(x,y)\to(0,0)}\frac{f(x,y)-x^2y^2}{(x^2+y^2)^2}=-1$。求 $f(0,0)$，判断原点极值类型，并证明存在 $\delta>0$，使 $0<x^2+y^2<\delta^2$ 时 $f(x,y)\le-\frac12(x^2+y^2)^2$。
\end{exercise}
\begin{exercise}{234}{三次多项式的驻点分类}
求 $z=x^2y-xy^2-xy$ 的全部驻点，逐一判定极值或鞍点并求极值；再判断函数在全平面是否有全局最大值或最小值。
\end{exercise}
\begin{exercise}{235}{指数因子的驻点与全局性质}
求 $f(x,y)=e^{-x}(x-y^3+3y)$ 的全部驻点及局部极值，分别给出二阶判别依据；再利用固定 $x$ 改变 $y$ 判断其全局有界性。
\end{exercise}
\begin{exercise}{236}{局部极值与全局界}
求 $f(x,y)=xe^{-(x^2+y^2)/2}$ 的全部驻点与极值，并证明所得极大值、极小值分别为全局最大值、最小值，写出全部取等点。
\end{exercise}
\begin{exercise}{237}{截面数据恢复与正定性}
设 $f(x,y)=2y^2+g(x-y)$，且对所有实数 $x$ 有 $f(x,0)=x^2$。先恢复 $g$ 与 $f$，再求极值，并分别用二阶判别及平方和表示证明其全局性质。
\end{exercise}
\begin{exercise}{238}{非孤立驻点与退化判别}
求 $f(x,y)=\frac12x^2+xy+\frac12y^2-2x-2y+5$ 的全部驻点和极值。说明为何二阶判别式为零仍可由配方得到结论，并区分严格极小与非严格极小。
\end{exercise}
\begin{exercise}{239}{隐函数的存在与极值}
设正值函数 $z=z(x,y)$ 由 $(x^2+y^2)z+\ln z+2(x+y+1)=0$ 确定。先证明每个 $(x,y)$ 对应唯一正根且隐函数处处可微，再求全部极值，判断所求极大值是否为全局最大值。
\end{exercise}
\begin{exercise}{240}{无盖水箱的全局最省材料}
设 $V>0$，用铁板制作容积为 $V$ 的无盖长方体水箱。求最省材料的长、宽、高及所需铁板面积，并证明所得驻点给出全局最小值，不能只列驻点方程。
\end{exercise}
\begin{exercise}{241}{闭区域内点与分段边界}
求 $f(x,y)=x^2+2y^2-x^2y^2$ 在 $D=\{(x,y):x^2+y^2\le4,\ y\ge0\}$ 上的最大值、最小值及全部取值点。要求分别列出内部驻点、直径边界与半圆边界的候选点并比较。
\end{exercise}
\begin{exercise}{242}{条件极值与偏导关系}
设 $f,\varphi\in C^1$，$P$ 为约束 $\varphi=0$ 下的局部极值点，且 $\varphi_y(P)\ne0$。证明 $f_x\varphi_y-f_y\varphi_x=0$。据此判断 $f_x(P)\ne0$ 是否必有 $f_y(P)\ne0$；再分别举例说明 $f_x(P)=0$ 时 $f_y(P)$ 可为零也可非零。
\end{exercise}
\begin{exercise}{243}{线性约束与距离最小化}
求 $f(x,y)=x^2+y^2-3$ 在 $x-y+1=0$ 下的全部条件极值，判断是否有最大值，并用原点到直线的距离解释最小值及其唯一性。
\end{exercise}
\begin{exercise}{244}{球内接长方体与全局最大体积}
设 $R>0$。求内接于半径为 $R$ 的球面的长方体的最大体积及长、宽、高。分别用拉格朗日乘数法与基本不等式证明结论，并说明退化边界不能取得最大值。
\end{exercise}
\begin{exercise}{245}{由梯度恢复函数后的区域最值}
设 $f(0,0)=1$，$f_x=6x$、$f_y=2y-1$。恢复 $f$，再求它在 $x^2+\frac{y^2}2\le1$ 上的最大值、最小值及全部取值点；需比较区域内部与椭圆边界。
\end{exercise}
\begin{exercise}{246}{相同二阶数据与不同极值类型}
构造三个二元多项式，使它们在原点的函数值、一阶偏导数和全部二阶偏导数完全相同，但原点分别为严格极小点、严格极大点和鞍点。逐一验证，并说明二阶判别式为零时不能仅凭这些数据决定极值类型。
\end{exercise}
\begin{exercise}{247}{四次型的参数分类与约束最值}
设 $F_a(x,y)=x^4+y^4+2ax^2y^2$，$a\in\mathbb R$。按参数完整分类原点的严格极小、非严格极小或鞍点性质；再求 $F_a$ 在单位圆 $x^2+y^2=1$ 上的最大值、最小值及全部取值点，单独讨论临界参数。
\end{exercise}
\subsection{二元函数的泰勒公式}
\begin{exercise}{248}{多项式泰勒展开与零余项}
设 $f(x,y)=2x^2-xy-y^2+3x+5y-1$。求它在 $(1,1)$ 处任意 $n\ge2$ 阶带拉格朗日余项的泰勒公式，并解释余项为什么恒为零；用展开式恢复原多项式作核验。
\end{exercise}
\begin{exercise}{249}{非原点的二阶泰勒公式}
求 $f(x,y)=x^3y^2$ 在 $(1,-1)$ 处带皮亚诺余项的二阶泰勒公式。令 $h=x-1,k=y+1$，明确写出常数、线性、混合二次项，并说明余项相对于 $h^2+k^2$ 的阶。
\end{exercise}
\begin{exercise}{250}{不同截断阶数的余项}
对 $f(x,y)=x^4+xy+(1+y)^2$，分别写出带皮亚诺余项的一阶、二阶麦克劳林公式，并求两种截断的精确余项，说明同一个 $x^4$ 项在两种余项条件下是否都可忽略。
\end{exercise}
\begin{exercise}{251}{三阶展开与二重极限}
求 $f(x,y)=\ln(1+x+y)$ 的三阶麦克劳林公式，余项取 $o((x^2+y^2)^{3/2})$；据此求 $\lim_{(x,y)\to(0,0)}\frac{\ln(1+x+y)+xy-x-y}{x^2+y^2}$，说明混合二次项如何抵消。
\end{exercise}
\begin{exercise}{252}{多元中值估计与符号区域}
设 $f$ 在全平面可微，$f_x>1$、$f_y<-1$、$f(0,0)=0$。分别讨论四个开象限中是否必有 $f(x,y)>x-y$，对必然成立或必然反向的情形作证明，对不能确定的情形给出两个线性函数作为例子。
\end{exercise}
\begin{exercise}{253}{二阶泰勒多项式的唯一性}
设两个次数不超过二的多项式 $P,Q$ 满足 $P(x,y)-Q(x,y)=o(x^2+y^2)$。证明 $P\equiv Q$。说明只沿两条坐标轴核验余项为什么不足以确定混合二次项，并给出明确的多项式反例。
\end{exercise}
\begin{exercise}{254}{退化二次项之后的极值判定}
设 $F_a(x,y)=e^{x+y}-1-x-y-a xy$。求原点处的三阶泰勒公式，并对全部实数 $a$ 判定原点是严格极小点、非严格极小点还是鞍点。二阶主项退化的参数须单独讨论，必要时沿直线检验更高阶项或使用指数函数不等式。
\end{exercise}
\subsubsection{综合提高题}
\begin{exercise}{255}{偏导符号与四点排序}
设 $f\in C^1(\mathbb R^2)$，$f_x>0$、$f_y<0$，记 $a=f(x,y)$、$b=f(x,y+1)$、$c=f(x+1,y)$、$d=f(x+1,y+1)$。确定四个数中必然的最大者与最小者并证明；再说明 $a,d$ 的次序能否确定，给出不同次序及相等情形的例子。
\end{exercise}
\begin{exercise}{256}{增长阶与可微性的四个命题}
设 $f$ 在原点连续。逐项判定下列命题，证明正确者并为错误者给出反例：
\begin{enumerate}
\item $\lim\frac{f(x,y)}{|x|+|y|}$ 为有限数，则 $f$ 在原点可微；
\item $\lim\frac{f(x,y)}{x^2+y^2}$ 为有限数，则 $f$ 在原点可微；
\item $f$ 在原点可微，则第一项中的极限存在；
\item $f$ 在原点可微，则第二项中的极限存在。
\end{enumerate}
所有极限均指 $(x,y)\to(0,0)$。再讨论第一项的极限值为零时可得到什么加强结论。
\end{exercise}
\begin{exercise}{257}{可微性与偏导数的定义区域}
设 $f(x,y)=(x^2+y^2)\sin\frac1{xy}$（$xy\ne0$），在 $xy=0$ 时取零。判断原点可微性并求全微分；在 $xy\ne0$ 处求 $f_x$，判断它能否连续延拓到原点；另用定义讨论坐标轴非原点处相应偏导数是否存在。
\end{exercise}
\begin{exercise}{258}{由二阶渐近判断极值与微分}
设 $\lim_{(x,y)\to(0,0)}\frac{f(x,y)-f(0,0)-x^3-2y^3}{\cos\sqrt{x^2+y^2}-1}=-2$。求 $f-f(0,0)$ 的二次主项，讨论原点连续性、两个偏导数、可微性及极值类型，并说明能否为极大点。
\end{exercise}
\begin{exercise}{259}{仅有两个偏导数时的结论边界}
已知 $f_x(0,0)=1$、$f_y(0,0)=2$。分别判断能否推出原点连续、存在全微分 $\dd x+2\,\dd y$、沿任意单位方向的方向导数为 $\cos\alpha+2\cos\beta$、沿 $x$ 轴负向的单侧方向导数为 $-1$。证明必然成立的结论，并构造反例说明其余结论不一定成立。
\end{exercise}
\begin{exercise}{260}{复合数据能否确定外函数偏导}
设 $f\in C^1((0,+\infty)^3)$，对所有 $x>y>0$ 有
\[
f(x-y,x+y,xy)=(x^2-y^2)^{xy}+(xy)^{2x}.
\]
判断这些条件能否唯一确定 $f_1(1,3,2)$，若不能，构造两个满足条件但该偏导数不同的函数。另令 $h(x,y)=f(x-y,x+y,xy)$，求可由已知数据确定的 $h_x(2,1)$，说明外函数偏导与复合函数导数的信息差别。
\end{exercise}
\begin{exercise}{261}{二阶偏导与截面初值}
设 $f_{yy}(x,y)=2$，$f(x,0)=1$、$f_y(x,0)=x$。求 $f$ 并证明在该数据下唯一，再判定原点是否为极值点。
\end{exercise}
\begin{exercise}{262}{沿曲线求导恢复二阶数据}
设 $f\in C^2$，$f_{xx}=f_{yy}$，且 $f(x,3x)=x$、$f_x(x,3x)=x^3$。求 $f_{xx}(x,3x)$ 和 $f_{xy}(x,3x)$。要求分别写出对两条给定关系求导得到的方程，再联立消元。
\end{exercise}
\begin{exercise}{263}{三元外函数的全微分}
设 $z=f(x,x+y,\frac xy)$，$f$ 可微，$y\ne0$。求 $\dd z$，并化简 $(x,y)=(0,1)$ 处的表达式，明确各偏导数在哪个三元点取值。
\end{exercise}
\begin{exercise}{264}{嵌套乘积的混合偏导}
设 $u=f(x,xy,xyz)$，其中 $f\in C^2$，$x,y,z$ 为独立变量。求 $u_{xy}$，化简 $u_{xy}(0,y,z)$，并说明不能先把 $x=0$ 代入使函数成为常数再求该混合偏导的原因。
\end{exercise}
\begin{exercise}{265}{零全微分的充要条件}
已知 $\lim_{(x,y)\to(0,0)}\frac{f(x,y)-f(0,0)}{\sqrt{x^2+y^2}}=0$。从定义求 $\dd f|_{(0,0)}$。再证明这一极限条件等价于原点可微且两个一阶偏导数均为零。
\end{exercise}
\begin{exercise}{266}{累次极限与联合趋近}
设 $x,y>0$，$f(x,y)=\frac y{1+xy}-\frac{1-y\sin(\pi x/y)}{\arctan x}$。先求 $g(x)=\lim_{y\to+\infty}f(x,y)$，再求 $\lim_{x\to0^+}g(x)$。判断联合趋近 $x\to0^+,y\to+\infty$ 时是否有相同极限，可比较 $y=x^{-2}$ 与 $y=x^{-3}$ 两条路径。
\end{exercise}
\begin{exercise}{267}{有理化与连续参数}
设 $a>0$，$f(0,0)=0$，非原点处 $f(x,y)=\frac{\sqrt{a+x^2y^2}-1}{x^2+y^2}$。讨论原点连续性的全部参数，并在连续的参数下进一步求两个偏导数及全微分；证明须覆盖任意趋近路径。
\end{exercise}
\begin{exercise}{268}{复合约束中的隐函数导数}
设 $u=f(x,y,z)$，$\varphi(x^2,y,z)=0$、$y=\sin x$，$f,\varphi\in C^1$ 且 $\varphi_3\ne0$。求 $\frac{\dd u}{\dd x}$。若 $\varphi(0,0,z_0)=0$，再写出经过该点的局部分支在 $x=0$ 的全导数，明确两组偏导数的取值点。
\end{exercise}
\begin{exercise}{269}{径向调和方程与初值}
设 $f\in C^2(0,+\infty)$，$z=f(\sqrt{x^2+y^2})$ 在原点外满足 $z_{xx}+z_{yy}=0$。推导 $f''(r)+\frac1r f'(r)=0$；若 $f(1)=0$、$f'(1)=1$，求 $f$，并讨论其原点处的有限连续延拓是否可能。
\end{exercise}
\begin{exercise}{270}{二阶方程的可逆变量变换}
设 $z\in C^2$ 满足 $z_{xx}+z_{xy}-z_{yy}=0$。求 $a,b$，使可逆变换 $u=x+ay,v=x+by$ 将方程化为 $z_{uv}=0$。计算变换后的非零比例系数，并说明为何不能选择 $a=b$。
\end{exercise}
\begin{exercise}{271}{隐式齐次关系与退化海森矩阵}
设 $z\ne0$，由 $\frac xz=\varphi(\frac yz)$ 确定 $C^2$ 局部分支，且 $x-y\varphi'(y/z)\ne0$。证明 $xz_x+yz_y=z$ 及 $z_{xx}z_{yy}=z_{xy}^2$；在第二个证明中说明为什么向量 $(x,y)$ 必非零，以及它与二阶偏导矩阵的关系。
\end{exercise}
\begin{exercise}{272}{空间交曲线的二阶隐式求导}
设 $a>0$，$z=x^2+y^2$、$x^2+2y^2+3z^2=4a^2$ 在关于 $(y,z)$ 的雅可比行列式非零的点附近确定 $y=y(x),z=z(x)$。求 $y'',z''$，最终表达式仅含 $x,y,z$；写出雅可比非退化条件，并利用 $z\ge0$ 将条件简化。
\end{exercise}
\begin{exercise}{273}{极值点数据与二阶复合导数}
设 $z=f(xy,yg(x))$，$f\in C^2$，$(1,1)$ 是 $f$ 的极值点，$g$ 由 $(x-1)^2-\int_x^{g(x)}e^{-t^2}\,\dd t=0$ 确定。求 $g(1),g'(1)$，并将 $z_{xy}(1,1)$ 用 $f$ 在 $(1,1)$ 的二阶偏导数表示。若该点为严格极小点，能否断言所得值严格为正？说明理由。
\end{exercise}
\begin{exercise}{274}{曲线数据确定切平面}
设 $f$ 可微，$f(t^2,t+1)=t^2(t-1)$，且 $f_1(1,2)=1$。求 $f(1,2)$、$f_2(1,2)$，再求曲面 $z=f(x,y)$ 在对应点的切平面与法线，说明曲线数据提供了哪个导数约束。
\end{exercise}
\begin{exercise}{275}{包含给定直线的椭球切平面}
在椭球面 $x^2+2y^2+3z^2=21$ 上求全部点 $M$，使过 $M$ 的切平面包含直线 $\frac{x-6}2=\frac{y-3}1=\frac{2z-1}{-2}$。求对应切平面，并把直线的基点条件与方向条件分别写出。
\end{exercise}
\begin{exercise}{276}{全微分相容性与最速增长}
设 $(3x^2+axy+2y^2)\,\dd x+(bxy+x^2)\,\dd y$ 是某个 $C^2$ 函数 $z$ 的全微分，且 $z(0,0)=1$。求 $a,b,z$，再求 $(1,0)$ 处方向导数的最大值及达到它的单位方向；用两个偏导数核验恢复结果。
\end{exercise}
\begin{exercise}{277}{二值隐函数的分支极值}
对方程 $x^2-6xy+10y^2-2yz-z^2+18=0$，先显式解出两个实值分支，证明它们在全平面正则。分别求两分支的全部极值点、极值，并判断相应极值是否为各分支的全局最值。
\end{exercise}
\begin{exercise}{278}{倒数约束与全局极值}
设 $x,y,z,a>0$，$\frac1x+\frac1y+\frac1z=\frac1a$。求 $f=xyz$ 的全部条件极值及取值点，证明其全局性质，并判断最大值是否存在，可将变量换为三个倒数。
\end{exercise}
\begin{exercise}{279}{双约束的可行参数区间}
在 $z=x^2+y^2$、$x+y+z=4$ 的约束下求 $u=x^2+y^2+z^2$ 的最大值、最小值及全部取值点。先推导可行的 $z$ 区间，再将目标化为 $z$ 的一元函数，验证区间端点确可实现。
\end{exercise}
\begin{exercise}{280}{约束曲面上的方向变化率极值}
在曲面 $2x^2+2y^2+z^2=1$ 上，求使 $f(x,y,z)=x^2+y^2+z^2$ 沿固定方向 $\vec l=(1,-1,0)$ 的方向导数最大的全部点及最大值；同时求最小值与取值点，并说明方向导数中的单位化。
\end{exercise}
\begin{exercise}{281}{椭球切面截体的最小体积}
设 $a,b,c>0$。在椭球面 $\frac{x^2}{a^2}+\frac{y^2}{b^2}+\frac{z^2}{c^2}=1$ 的 $x,y,z>0$ 部分求点 $P$，使其切平面与三个坐标面围成的四面体体积最小。求点、切平面和最小体积，并解释趋向坐标边界不能产生更小体积。
\end{exercise}
\begin{exercise}{282}{切平面投影域上的最值}
设 $f(x,y)=x^3+y^3-(x+y)^2+3$。求曲面 $z=f(x,y)$ 在 $(1,1,1)$ 处的切平面 $T$，并写出 $T$ 与三个坐标面围成的有界区域在 $xOy$ 面上的投影 $D$。再求 $f$ 在 $D$ 上的最大值、最小值及全部取值点，分别检查三条边与内部。
\end{exercise}
\begin{exercise}{283}{由边际成本恢复总成本与约束优化}
某企业生产甲、乙两种产品的固定成本为 $10000$ 万元，产量为 $x,y$ 件，边际成本分别为 $C_x=20+\frac12x$、$C_y=6+y$，单位为万元/件。求总成本 $C(x,y)$。在 $x+y=50$、$x,y\ge0$ 时求最省成本的产量及最小成本；先按实数产量求解，再判断非负整数产量限制是否改变最优结果。
\end{exercise}
